% Generated draft: requires mathematical and source review.
\documentclass[11pt]{article}
\usepackage[margin=25mm]{geometry}
\usepackage{fontspec}
\setmainfont{Latin Modern Roman}
\usepackage{amsmath,amssymb,mathtools,unicode-math}
\setmathfont{Latin Modern Math}
\usepackage{enumitem,needspace,xurl,hyperref}
\setlength{\parindent}{0pt}
\setlength{\parskip}{5pt}
\setlength{\emergencystretch}{3em}
\newcommand{\N}{\mathbb{N}}
\newcommand{\Z}{\mathbb{Z}}
\newcommand{\Q}{\mathbb{Q}}
\newcommand{\R}{\mathbb{R}}
\newcommand{\C}{\mathbb{C}}
\newcommand{\F}{\mathbb{F}}
\newcommand{\A}{\mathbb{A}}
\newcommand{\PP}{\mathbb{P}}
\newcommand{\E}{\mathbb{E}}
\newcommand{\eps}{\varepsilon}
\newcommand{\dd}{\,\mathrm{d}}
\begin{document}
\section*{UnsolvedMath ID: 30001223}
\textbf{Source problem code:} catalogue entry
\subsection{Definitions and mathematical statement}
A partition \(\lambda\) of an integer \(r\geq 0\) is a finite weakly decreasing sequence of positive integers whose sum is \(r\); the empty partition has sum zero. Write \(\ell(\lambda)\) for the number of parts. The source does not specify the range of \(r\) or whether the empty partition is included.

For a finite-dimensional module \(M\), define its radical \(\operatorname{rad}M\) to be the intersection of its maximal proper submodules, and its top by
\[
\operatorname{top}M=M/\operatorname{rad}M.
\]
A module is simple if it is nonzero and has no nonzero proper submodule.

Retain the source notation \(Y^\lambda\) for the Young module and \(I_n(\lambda)\) for the indecomposable injective polynomial module indexed by \(\lambda\), with \(n\geq\ell(\lambda)\). Retain \(\Lambda^{\mathrm{in}}\) for the De Visscher--Donkin index set. The supplied record does not specify the ground field, its characteristic, the ambient polynomial representation category, or the definition of this index set. In particular, injectivity and projectivity must be interpreted in the intended common category, which has not been supplied.

\textbf{Proposed equivalence, subject to these unresolved specifications.} For every partition \(\lambda\) of \(r\) in the intended setting, is it true that
\[
\begin{aligned}
&\operatorname{top}Y^\lambda\text{ is simple}\\
&\quad\Longleftrightarrow\quad
\exists n\in\Z_{\geq 1}\;\bigl(n\geq\ell(\lambda)\ \text{and}\ I_n(\lambda)\text{ is projective}\bigr)\\
&\quad\Longleftrightarrow\quad
\lambda\in\Lambda^{\mathrm{in}}\,?
\end{aligned}
\]
Here positive integral \(n\) is the conventional interpretation of the polynomial-module rank parameter; the source gives only the inequality \(n\geq\ell(\lambda)\). Both displayed equivalences are retained from the source. The index-set equivalence cannot be made independently precise without its definition, and the displayed question requires the missing representation-theoretic hypotheses before it constitutes a fully specified conjecture.
\subsection{Short English statement}
The supplied conjecture asks whether a Young module has simple top exactly when its partition labels a projective indecomposable injective polynomial module for some sufficiently large rank, and equivalently exactly when the partition lies in the De Visscher--Donkin set \(\Lambda^{\mathrm{in}}\). The ground field, characteristic, module category, and index-set definition are missing. The supplied open status is dated curation, not an independently verified conclusion.
\subsection{Sources}
\begin{itemize}
\item Local UnsolvedMath registry, supplied complete source record: permanent ID \(30001223\), source problem code \texttt{catalogue entry}, title \emph{Simple Tops of Young Modules}, in the Oberwolfach Reports Open Problems collection.
\item Supplied report citation: \emph{Representations of Finite Groups} (2009), Oberwolfach Reports, pp.~917--970, DOI 10.4171/owr/2009/17. Source URL and literature source labelled ``original report'': \url{https://doi.org/10.4171/owr/2009/17}. The supplied background describes these reports as documenting mathematical workshops and their open-problem discussions.
\item Supplied literature source labelled ``De Visscher Donkin'': \url{https://openaccess.city.ac.uk/id/eprint/969/1/prinjjan12.pdf}. The record associates the classification conjecture and index set with De Visscher--Donkin, but supplies no full bibliographic entry or individual proposer.
\item Literature assessment checked at \texttt{2026-08-21}: the supplied curation labels the problem open and describes the extracted De Visscher--Donkin classification conjecture as an unresolved general problem about simple tops of Young modules. This is supplied dated curation derived from the OWR material, not an independently verified assessment; the source itself calls for independent verification of time-sensitive status and cited claims.
\item The supplied statement-status label is \texttt{corrected\_verified}. Its accompanying verification concerns balanced mathematical delimiters and retention of both equivalences; it does not establish the missing hypotheses or mathematical correctness. No cited document was independently consulted for this draft.
\end{itemize}
\end{document}
